\chapter{Alain Connes}

\addcontentsline{toc}{chapter}{Problem History}
\section*{Problem History}

\begin{prereqs}
\noindent \textbf{Background}
\begin{itemize}
\item Functional analysis and operator algebras
\item Topology of manifolds and differential geometry
\item Measure theory and ergodic theory
\end{itemize}
\end{prereqs}

\begin{objectives}
\noindent \textbf{Learning Objectives}
\begin{itemize}
\item Understand the classification program for von Neumann algebras and the role of factors
\item Learn the concepts of noncommutative geometry: spectra of noncommutative algebras and cyclic cohomology
\item Appreciate how Connes unified index theory, foliation theory, and operator algebras into a coherent framework
\end{itemize}
\end{objectives}

Von Neumann algebras---weakly closed self-adjoint subalgebras of bounded
operators on a Hilbert space---sit at the crossroads of operator theory,
ergodic theory, and geometry. Murray and von Neumann classified factors
into three types: Type $I$ (bounded operators on a Hilbert space), Type
$II_1$ (unique faithful normal tracial state), Type $II_\infty$, and
Type $III$ (arising in ergodic theory and quantum field theory).

Connes's classification \cite{connes1985} of type $III$ factors (1972--1973) introduced the
flow of weights $T(M)$, which takes values in positive reals modulo
multiplication. This invariant, defined through Takesaki's modular
automorphism groups, gives a complete classification: two type $III$
factors are isomorphic iff their flows of weights are isomorphic. Every
type $III$ factor is a crossed product of a commutative algebra by a
measurable flow, connecting operator algebras with dynamical systems.

Connes then created noncommutative geometry. The key insight: for a
compact space $X$, the commutative $C^*$-algebra $C(X)$ determines $X$.
One should study ``noncommutative spaces'' via their noncommutative
algebras. Concrete tools---particularly cyclic cohomology---enabled
genuine geometry on such spaces.

\addcontentsline{toc}{chapter}{How Was It Solved}
\section*{How Was It Solved}

The Hochschild cohomology $HH^*(A)$ of an algebra $A$ plays a role
analogous to de Rham cohomology. Connes introduced cyclic cohomology
$HC^*(A)$ as a refinement, consisting of cochains invariant under cyclic
permutation. For a smooth manifold $M$, $HC^*(C^\infty(M))$ is isomorphic
to de Rham cohomology, and the pairing with $K$-theory generalizes the
Chern character.

The fundamental theorem is the index theorem for noncommutative spaces.
For a spectral triple $(A, H, D)$---an algebra $A$ on Hilbert space $H$
with Dirac operator $D$---the index formula extends as a pairing between
a cyclic cocycle and the $K$-theory of $A$.

The local index formula (Connes--Moscovici) expresses the index via the
Dixmier trace---a singular trace on the Macaev ideal
$\mathcal{L}^{1,\infty}$ providing a noncommutative integral. The index
becomes a local heat-kernel expression analogous to McKean--Singer.

Foliations provided the testing ground. The holonomy groupoid of a
foliation is a noncommutative groupoid; its $C^*$-algebra encodes the
geometry. The index theorem for twisted Dirac operators involves the
Godbillon-Vey invariant.

Spectral triples (1990s) axiomatize noncommutative Riemannian geometry:
$A$ represented on $H$, $D$ with compact resolvent, bounded commutators
$[D, a]$ for $a \in A$.

\addcontentsline{toc}{chapter}{Impact}
\section*{Impact}

The flow of weights transformed von Neumann algebra theory from
classification results to deep structural insights. Noncommutative
geometry has found applications across mathematics and physics.

In number theory, Connes formulated a noncommutative approach to the
Riemann hypothesis using the adèle class space
$\mathbb{A}_\mathbb{Q}/\mathbb{Q}^*$. In mathematical physics, Connes
and Chamseddine showed the Standard Model arises from a spectral triple:
the algebra encodes $SU(3) \times SU(2) \times U(1)$ and the Higgs
field appears as a scalar field. This has been extended to incorporate
gravity.

Cyclic cohomology is essential in geometry and topology, with the local
index formula generalized to manifolds with boundary and singular spaces.

\addcontentsline{toc}{chapter}{Modern Views}
\section*{Modern Views}

The reconstruction theorem for commutative $C^*$-algebras led to
``almost-commutative'' geometries unifying the Standard Model with
general relativity. Extended Connes-Chamseddine models accommodate the
observed Higgs mass of 125 GeV.

In number theory, the noncommutative trace formula encodes Riemann zeta
zeros, connecting to quantum chaos. Cyclic cohomology has found
applications in topological quantum field theory and topological phases
of matter; the noncommutative torus is fundamental in quantum Hall effect
theory.

Connes received the Fields Medal in 1982 for operator algebras and
noncommutative geometry.

\section*{Summary}
Connes classified type $III$ von Neumann algebras using the flow of
weights, then created noncommutative geometry with cyclic cohomology and
spectral triples. The framework has applications to number theory,
particle physics, and quantum gravity, and continues to evolve as a
unifying perspective across mathematics and physics.

\addcontentsline{toc}{chapter}{Exercises}
